\section{Performance Landscape Analysis}
\label{sec:app_landscape}

This section reasons through the expected performance landscape for each
LLAMIA-Bench task, identifying what information each system type can and
cannot access and predicting where the verb--latent gap, the SFT--DAPO gap,
and the interface$\times$optimization interaction should be largest.
The analysis is grounded in the task structure and the literature on each
capability; it serves as a prior against which empirical results can be
interpreted.

Throughout, we consider six system types in a consistent order:

\begin{enumerate}[leftmargin=*, itemsep=1pt, topsep=2pt]
  \item \textbf{LLM Only}: Qwen3 backbone with no subagent access (FEN input, pretraining knowledge only).
  \item \textbf{Frontier\,+\,Lc0 Verb}: GPT-5 with Lc0-BT4 via ReAct tool-calling (top-$k$ moves, evals, WDL, PVs as text).
  \item \textbf{LLAMIA-Verb-SFT}: Verbalized subagent output, single-step SFT.
  \item \textbf{LLAMIA-Verb}: Verbalized subagent output, DAPO.
  \item \textbf{LLAMIA-SFT}: Latent tokens ($k{=}32$), single-step SFT.
  \item \textbf{LLAMIA}: Latent tokens ($k{=}32$), DAPO.
\end{enumerate}

\noindent Dedicated task experts (Allie, Chessformer, SCC, etc.) are discussed
where applicable.

\paragraph{Two sources of verbalization debt.}
Across all tasks, the verb--latent gap arises from two interacting sources.

The first is \emph{non-verbalizable information}.
The engine's internal representation encodes structures that have no faithful
text equivalent: learned lookahead---compressed representations of tactical
sequences three to eight moves deep, including threats that never materialize
on the board because the opponent avoids them; attack geometry---which
diagonals bear pressure, which squares around the king are weak, how piece
coordination creates threats; strategic themes---that a pawn recapture
permanently damages the dark-square complex, that an isolated pawn will become
an endgame target, that the bishop pair grows in strength as the position
opens; and distributional properties of the policy---which alternatives are
genuinely competitive, how sharply the evaluation landscape changes across
candidate moves.
Verbalization collapses this representational state into a scalar evaluation,
a ranked move list, and a principal variation in algebraic notation---a lossy
projection that preserves \emph{what} and \emph{how much} but destroys
\emph{why} and \emph{what distinguishes}.

The second is \emph{collaboration agency}.
Access to the non-verbalizable information enables productive collaboration
strategies---counterfactual queries, consult-then-override, multi-step
lookahead---that are structurally unproductive under verbalization
(\cref{sec:abl_agency_strategies}).
The debt is not merely perceptual (richer information per invocation) but
agential (richer strategies for \emph{using} that information).
The two sources interact: the information enables the agency, and the agency
amplifies the information advantage across invocations.

The per-task analyses below identify which non-verbalizable features and
which collaboration strategies are load-bearing for each task, explaining
both the magnitude and the ordering of the verb--latent gap across
LLAMIA-Bench.

% =============================================================================
\subsection{Puzzle Interest Estimation}
\label{sec:app_landscape_interest}
% =============================================================================

\paragraph{What the task requires.}
The model predicts whether a puzzle position is \emph{interesting to a human
player}: a subjective, aesthetic judgment that depends on sacrifices, quiet
winning moves hiding deep tactics, stalemate traps, positions with
counter-intuitive evaluations, and the presence of multiple competing
strategic plans.
No prior dedicated model exists for this task; LLAMIA-Bench introduces it as
a new evaluation target.

\paragraph{The information asymmetry.}
Interest is fundamentally about the \emph{shape} of the engine's internal
distribution, not any single number the engine outputs.
A position is interesting when the engine's policy concentrates on a
counter-intuitive move (sacrifice, quiet retreat), when the value function
has a large gradient across similar-looking moves (one move wins, the obvious
one loses), or when the distribution has unusual structure compared to
typical positions.
Verbalization collapses this distributional object into a handful of
numbers---top-$k$ moves and their centipawn evaluations---destroying exactly
the signal that defines interest.

\paragraph{Per-system analysis.}

\emph{LLM Only.}
The model knows from pretraining that sacrifices and stalemate traps are
interesting, but without positional evaluation it cannot \emph{detect} these
features in a given FEN\@.
Predictions reduce to surface heuristics (piece count, game phase), yielding
near-random correlation.

\emph{Frontier\,+\,Lc0 Verb.}
The verbalized output surfaces some signals: a large evaluation gap between
the top move and the second-best suggests a forcing combination; a PV
containing material imbalance suggests a sacrifice.
However, these heuristics cover only a fraction of what makes a position
interesting.
The frontier model falls back on piece count and game-phase reasoning,
achieving only marginally above the residual guess.

\emph{Verbalized baselines (SFT and DAPO).}
Through training, verbalized baselines learn the same correlations the
frontier model discovers (evaluation gaps, PV features).
However, the mapping from verbalized features to interest has high
irreducible noise because the distributional information that defines
interest is absent.
Both verb baselines regress toward predicting the dataset mean, reaching
a local optimum in which a near-constant prediction minimizes training loss.
The SFT--DAPO gap is small because interest estimation is a single-step
perceptual judgment with no multi-step reasoning for RL to improve.

\emph{LLAMIA-SFT (latent).}
The latent tokens encode the engine's full policy distribution and
positional feature activations---exactly the signal that distinguishes
interesting from routine positions.
A position where the engine's policy concentrates on a sacrifice has a
characteristically different activation pattern from one with a standard
best move.
SFT can learn this mapping directly, yielding a substantial jump over
verbalized baselines.

\emph{LLAMIA (latent\,+\,DAPO).}
Interest is a single-step judgment: one position, one scalar prediction.
RL's comparative advantage (learning to sequence actions across steps) does
not activate.
The marginal gain from DAPO comes only from sharpening the latent-to-interest
mapping under the correlation reward, which is noisy.
The SFT--DAPO gap should therefore be the \emph{smallest} of any LLAMIA-Bench
task.

\paragraph{Expected signatures.}
Verbalized baselines should produce near-constant predictions (low output
variance); LLAMIA should produce high-variance predictions that correlate
with actual interest.
If predicted-vs.-actual interest is plotted, verbalized baselines form a
horizontal band while LLAMIA shows a positive trend.
The verb--latent gap on interest should be the \emph{largest} among all tasks
(measured as relative improvement), because interest is the task where
verbalization loses the most decision-relevant information.

\paragraph{Predicted behavior summary.}
\begin{itemize}[leftmargin=*, itemsep=2pt, topsep=2pt]
\item \textbf{Scale (4B${\to}$14B):} Large under latent
  ($3.1{\times}$ in \cref{fig:scaling_bars}); negligible under
  verbalized---the discriminative signal is absent from the verbal
  channel regardless of backbone capacity.
\item \textbf{SFT vs.\ DAPO:} \emph{Smallest gap in the benchmark}.
  Single-step perceptual judgment with no sequencing for RL to improve;
  DAPO gain comes only from sharpening the latent-to-interest mapping
  under a noisy correlation reward.
\item \textbf{Verbalization debt:} \emph{Largest gap in the benchmark}.
  Policy shape, value gradients, and structural asymmetries have no
  verbal proxy; no collaboration strategy under verbalization can
  recover them.
\item \textbf{vs.\ Frontier\,+\,Lc0:} Marginally above residual guess.
  Inference capacity cannot recover absent signal; the model falls back
  on piece count and game-phase heuristics.
\item \textbf{vs.\ task experts:} None---LLAMIA-Bench introduces
  Interest as a new evaluation target.
\end{itemize}

% =============================================================================
\subsection{Puzzle Difficulty Estimation}
\label{sec:app_landscape_difficulty}
% =============================================================================

\paragraph{What the task requires.}
The model predicts how hard a puzzle is for human solvers.
Difficulty depends on: solution depth, whether the first move is forcing
(checks, captures) or quiet, the number of plausible-looking distractors,
required calculation depth, and pattern familiarity.

\paragraph{Why the verb--latent gap is smaller than Interest.}
Unlike interest, difficulty has a \emph{moderately informative verbal
proxy}: solution length.
The principal variation length is directly readable from the verbalized engine
output, and solution length correlates substantially with human-perceived
difficulty.
Verbalization therefore loses less decision-relevant information for
difficulty than for interest.

\paragraph{Per-system analysis.}

\emph{Frontier\,+\,Lc0 Verb.}
The frontier model reads the PV length and adds heuristic adjustments:
endgames are harder (more technique required), quiet first moves are harder
than forcing ones.
This first-order approximation achieves moderate rank correlation.

\emph{Verbalized baselines (SFT and DAPO).}
Training learns the PV-length correlation strongly.
May additionally learn secondary features: the evaluation gap between the
best move and the second-best as a proxy for ``distractor strength,'' and
game-phase indicators.
The ceiling is reached when two puzzles share the same solution length but
differ in cognitive difficulty (e.g., a three-move forcing sequence of checks
vs.\ a three-move quiet combination).
The verbalized signal cannot distinguish these.

\emph{LLAMIA-SFT (latent).}
The engine's activations encode the \emph{nature} of the required
computation, not merely its length.
A position requiring deep quiet play produces different activation patterns
from one with a forced tactical sequence.
The policy distribution shape is informative: a sharply peaked distribution
indicates an ``obvious'' first move (easier), while a flat distribution
indicates genuine ambiguity (harder).
Value gradients across candidate moves indicate how much calculation is
needed to distinguish them.
Concretely, latent reads recover patterns that humans recognize at a
glance but that have no clean text proxy: back-rank mate motifs,
trapped-piece geometries, forced-move sequences (each reply is the only
legal response, collapsing search width), and ``mate-in-$N$ where the
moves are obvious'' positions that share solution length with genuinely
hard combinations but require far less calculation.
These cues compress difficulty along an axis orthogonal to PV length,
which is what the verb--latent residual on this task captures.

\emph{LLAMIA (latent\,+\,DAPO).}
RL helps more on difficulty than on interest.
Difficulty estimation benefits from the model learning to ``simulate'' the
solving process: how many steps of reasoning are required, and how
uncertain is the intermediate evaluation at each step?
DAPO rollouts expose the model to sequential position evaluation, which
calibrates difficulty estimates.
The SFT--DAPO gap should be slightly larger than for interest, but
substantially smaller than for commentary or planning tasks.

\paragraph{Expected signatures.}
Both verbalized and latent systems should produce \emph{correlated}
predictions, but the latent system should discriminate better among puzzles
with equal solution length.
A residual analysis conditioning on PV length should reveal that verbalized
baselines explain negligible additional variance beyond length, while LLAMIA
captures computational-complexity features orthogonal to length.

\paragraph{Predicted behavior summary.}
\begin{itemize}[leftmargin=*, itemsep=2pt, topsep=2pt]
\item \textbf{Scale (4B${\to}$14B):} Small-to-moderate.
  PV length proxy ceilings verb baselines regardless of LLM capacity;
  latent residual is narrow enough that larger backbones add only
  incremental gain.
\item \textbf{SFT vs.\ DAPO:} Slightly larger than Interest,
  substantially smaller than Annotation or Commentary.
  DAPO rollouts expose the model to sequential position evaluation,
  which calibrates difficulty estimates beyond single-example SFT.
\item \textbf{Verbalization debt:} Moderate, slightly smaller than BC.
  Solution depth absorbs most variance; the latent residual captures
  computational-character features (forced-move sequences, back-rank
  motifs, distractor sharpness) orthogonal to length.
\item \textbf{vs.\ Frontier\,+\,Lc0:} Moderate rank correlation via
  PV length plus heuristic adjustments (endgames harder, quiet first
  moves harder than forcing ones).
\item \textbf{vs.\ task experts:} None; closest reference is human
  solve-rate calibration (Lichess puzzle rating), which uses
  population-aggregate statistics rather than position-conditioned
  prediction.
\end{itemize}

% =============================================================================
\subsection{Behavior Cloning}
\label{sec:app_landscape_bc}
% =============================================================================

\paragraph{What the task requires.}
Given a position and a target Elo, predict the move a human at that skill
level would play.
The metric is move-matching accuracy.
LLAMIA-Bench evaluates two splits: \emph{MAIA} (standard Elo-stratified
buckets) and \emph{Wild} (GM-25 opening moves, Low-Time blitz positions,
Elo-Gap games where players deviate from their rating).

\paragraph{What the literature establishes.}
\citet{mcilroy2022aligning} demonstrated that engine strength and human move
prediction are orthogonal: Stockfish depth-15 achieves 33--41\% move-matching
accuracy (monotonically increasing with target Elo, peaking on the 1900
bucket), while Maia's per-Elo CNNs reach ${\sim}$52\% with unimodal accuracy
curves that peak near each model's training Elo.
Leela Zero caps at ${\sim}$46\% on the 1900 bucket and is flat across lower
buckets, confirming that raw engine strength does not transfer to human
prediction.
\citet{zhang2024humanalignedchessbitsearch} (ALLIE) pushed this to 55.7\%
(policy head alone) with a unified 355M-parameter decoder transformer and
soft Elo conditioning that linearly interpolates between weak- and
strong-player embeddings; adaptive MCTS at inference further achieves
49~Elo MAE skill calibration.
Critically, ALLIE's ablations show the model is \emph{data-constrained, not
parameter-constrained}: halving parameters drops accuracy by only 1.2~pp,
doubling them adds only 0.3~pp, while halving training data costs 1.0~pp.
MAIA-3~\cite{XXX_chessformer} achieved 57.1\% at only 79M parameters through
domain-aligned architecture (square tokenization, geometric attention bias,
source-destination policy head), surpassing search-enabled ALLIE at
$4.5{\times}$ fewer parameters---and GPT-3.5 (53.7\% at 175B) at
$2000{\times}$ fewer parameters.
The ${\sim}$57\% ceiling likely reflects genuine stochasticity in human move
selection: in most positions multiple reasonable moves exist and the
``correct'' human move depends on unobservable factors (time pressure,
preparation, psychological state).

\paragraph{Key factors from the literature.}
Three factors dominate BC performance, in order of impact:
(i)~\emph{domain-aligned architecture}---Chessformer's inductive biases
(square tokens, geometric attention, from-to policy head) yield a larger
gain than $2000{\times}$ more parameters in a generic architecture,
establishing that integration interface dominates backbone scale;
(ii)~\emph{skill conditioning}---continuous Elo conditioning (ALLIE's soft
tokens) outperforms discrete binning (Maia's per-Elo models) by enabling
interpolation and shared representations across the skill spectrum;
(iii)~\emph{training objective}---RL over imitation
(\cite{XXX_c1} shows +7~pp from SFT to DAPO on chess puzzles at 4B scale,
demonstrating that RL provides gains beyond maximum-likelihood imitation,
though the effect may be smaller on BC where the reward signal is noisier
than puzzle correctness).

\paragraph{Per-system analysis.}

\emph{LLM Only.}
Qwen3 has seen chess games in pretraining but lacks board-state tracking
fidelity~\cite{XXX_chessqa}.
\cite{XXX_c1} confirms that untrained Qwen3 at 4B scale produces near-zero
accuracy on chess puzzles, and open-source models ${\leq}$14B without
chess-specific distillation remain in the single-to-low-double digits.
Without subagent access or task-specific training, the model produces
``literary'' moves---plays it has encountered in chess commentary and
book annotations.
It has no Elo-conditioning mechanism and cannot distinguish 1200-level play
from 2000-level play.
Performance should be well below the Maia baseline, scaling with model size
as more chess pretraining knowledge becomes available at 8B and 14B.

\emph{Frontier\,+\,Lc0 Verb.}
The engine's top-$k$ are all strong moves, but humans below ${\sim}1800$ Elo
frequently play moves not in the engine's top-5.
The frontier model must ``de-tune'' engine suggestions to match the target
Elo---a task for which verbalized evaluations provide only coarse guidance.
The model can apply heuristics (``at 1200, prefer the simpler-looking
move'') but cannot access the distributional structure that indicates which
alternative moves the engine considers ``close'' in its representation space.
A lower bound is the engine pass-through ceiling: if the model simply copies
Lc0's top move, it recovers Leela's ${\sim}$46\% on the 1900 bucket but
collapses at lower Elos where engine-best diverges most from human play.
Performance should exceed LLM-Only substantially but plateau below dedicated
BC experts.

\emph{Verbalized baselines (SFT and DAPO).}
Training on the actual data distribution teaches the model which patterns in
verbalized engine output correlate with human move choice at each Elo.
The model learns that when the evaluation gap between move~1 and move~3 is
small and move~3 involves a simpler piece interaction, lower-Elo players
prefer move~3.
However, the top-$k$ moves cover perhaps 60--70\% of what humans actually
play; moves outside this set are invisible to the verbalized system.
RL adds moderate value by directly optimizing move-matching accuracy rather
than cross-entropy, but the information ceiling is the same.

\emph{LLAMIA-SFT (latent).}
The full policy distribution over all legal moves is present in the latent
tokens---not just the top-5 but the shape of the distribution.
For BC this matters: the distribution encodes which moves the engine
considers ``similar'' in representation space.
At low Elo, humans tend to play moves that are similar to good moves in the
engine's feature space (conceptually nearby but tactically inferior), and
this similarity structure is exactly what verbalization destroys.
The latent tokens also encode positional features (king safety, pawn
structure, piece activity) whose \emph{relative importance changes with
skill level}: strong players weigh king safety more heavily while weaker
players fixate on material.
SFT learns the mapping from (latent state, Elo) $\to$ human move on
single examples.

\emph{LLAMIA (latent\,+\,DAPO).}
RL adds meaningful value on BC because the reward signal (move match) is
clear and non-noisy.
DAPO can discover non-obvious mappings between latent features and human
moves that SFT's maximum-likelihood objective misses---in particular, it
can learn that for certain position types, the best prediction is the
\emph{second} peak of the bimodal policy distribution, not the primary
peak.
Performance may approach but is unlikely to exceed the Chessformer ceiling
(${\sim}$57\%): LLAMIA's adapter is shallow (3-layer MLP), its Stage-2
training is multi-task (not BC-specific), and it lacks the domain-aligned
inductive biases (square tokenization, geometric attention) that proved
decisive in the Maia-3 results.

\emph{Dedicated experts (Allie, Chessformer).}
End-to-end optimization for exactly this objective, with domain-aligned
architecture and vastly more training data (93M games for Allie).
These represent the performance ceiling on standard MAIA buckets.
ALLIE's published 55.9\% on Lichess blitz~\cite{zhang2024humanalignedchessbitsearch}
and Chessformer's 57.1\%~\cite{XXX_chessformer} bound the in-distribution
BC expert range.

\paragraph{MAIA vs.\ Wild splits.}
On standard MAIA buckets, dedicated experts should retain a slight edge
because their architectural specialization is well-matched to the
move-prediction objective.
On Wild splits, the landscape shifts because each split introduces a distinct
out-of-distribution challenge:

\emph{GM-25.}
LLAMIA-Bench provides ${\leq}$4.6K games per GM---well below the
${\sim}$150K threshold that Maia-Individual~\cite{XXX_maia_individual}
identifies as necessary for meaningful per-player adaptation.
Allie and Chessformer, trained on aggregate Lichess blitz, collapse to their
population predictor in this regime; elite GMs' primary games are not played
on Lichess blitz, making this split doubly OOD (style + time control).
LLAMIA's latent tokens may encode position-evaluation features that
correlate with GM-specific style preferences, but the expected gain over
verbal baselines is modest because per-player adaptation requires signal
the engine does not carry.
The LLM's pretraining knowledge of opening theory, named players, and
their characteristic systems provides a partial advantage on this split
that is orthogonal to both the latent channel and the BC experts'
position-conditioned training; this is an LLM-backbone effect, not a
latent-internalization effect, and the verb--latent gap on GM-25 should
be close to constant across the Elo-bucket Lichess data.

\emph{Low-Time.}
Maia's original protocol discards all moves with ${<}$30\,s on the
clock~\cite{mcilroy2022aligning}; ALLIE's training includes time pressure
implicitly through Lichess blitz but not as explicit conditioning metadata.
BT4's latent state does not encode clock information, so latent injection
does \emph{not} provide a structural advantage here---the bottleneck is
missing signal, not interface.
Any Low-Time advantage LLAMIA shows over LLAMIA-Verb at fixed clock
metadata is therefore an LLM-backbone effect---pretraining knowledge of
time-management heuristics such as simplification under flag and a bias
toward principled moves when short on time---not evidence that the
latent channel carries clock information.
Dedicated BC experts with blitz-heavy training data (ALLIE) retain most
of their in-distribution accuracy.

\emph{Elo-Gap.}
ALLIE explicitly conditions on both player ratings as soft
tokens~\cite{zhang2024humanalignedchessbitsearch}, so mismatched-strength
games are inside its conditioning grammar, though underrepresented in
training data.
LLAMIA's latent tokens carry the full policy distribution, which can
distinguish positions where the stronger player offers tactical
complications from positions where the weaker player seeks simplification.
This positional asymmetry is present in BT4's activations but not cleanly
verbalizable; the latent advantage should be moderate.

\paragraph{Expected qualitative behaviors and accuracy curves.}
The per-Elo accuracy curve is a diagnostic:
engines (Stockfish, Leela) show \emph{monotonically increasing} accuracy
with target Elo (best at 1900, worst at 1100) because engine-best moves
align with strong human play;
Maia-style models show \emph{unimodal} curves peaking near their training
Elo~\cite{mcilroy2022aligning};
ALLIE and Chessformer show \emph{flat-to-slightly-decreasing} curves across
Elos because they condition on skill level.

\emph{LLM Only}: should mirror the engine-like monotonically increasing
curve---produces book moves gravitating toward strong play, with no Elo
sensitivity.
\emph{Verb baselines}: DAPO Elo-conditioning should bend the curve toward
Maia-style unimodal behavior, but verbalized PV cannot carry the
skill-style features that ALLIE's soft Elo tokens encode; expect a flatter
but still engine-skewed curve.
Over-predicts engine-optimal moves for low-Elo targets---plays
``engine-lite'' rather than characteristically human.
Cannot reproduce skill-specific error patterns (e.g., the systematic
material overvaluation of 1200-rated players).
\emph{LLAMIA}: closest to ALLIE's flat profile because latent state carries
skill-conditioning directly.
Should reproduce Elo-specific behavioral signatures: does LLAMIA-1200 fall
for opening traps that Maia-1200 falls for?
Does LLAMIA-2000 exhibit the material indifference characteristic of strong
positional players?
\emph{Dedicated experts}: most faithful move-matching at each Elo, including
error-pattern reproduction, but mute---cannot explain their predictions or
transfer to adjacent tasks.

\paragraph{$2{\times}2$ factorial prediction.}
The interface and optimization effects should be roughly additive on BC\@.
BC is a single-step task (predict one move per position), so the multi-step
advantage of DAPO is limited; the RL gain is small under both
representations, with no strong interaction effect.
This contrasts with Commentary, where the interaction is super-additive.

\paragraph{Predicted behavior summary.}
\begin{itemize}[leftmargin=*, itemsep=2pt, topsep=2pt]
\item \textbf{Scale (4B${\to}$14B):} Smallest relative effect among
  representation-intensive tasks ($2.4{\times}$ in
  \cref{fig:scaling_bars}); ALLIE's ablations establish that BC is
  data-constrained, not parameter-constrained.
\item \textbf{SFT vs.\ DAPO:} Small gap, additive
  interface${\times}$optimization.
  Single-step prediction with a clear non-noisy reward (move match);
  RL sharpens without discovering new collaboration strategies.
\item \textbf{Verbalization debt:} Moderate, slightly larger than
  Difficulty.
  Top-$k$ covers ${\sim}$60--70\% of human moves; the residual
  distributional similarity structure carries skill conditioning and
  is accessible only to the latent channel.
\item \textbf{vs.\ Frontier\,+\,Lc0:} Above LLM-Only substantially;
  plateaus below dedicated BC experts.
  Engine pass-through recovers Leela's ${\sim}$46\% on the 1900 bucket
  and collapses at lower Elos.
\item \textbf{vs.\ task experts:} ALLIE (55.9\%) and Chessformer
  (57.1\%) bound the in-distribution MAIA ceiling.
  LLAMIA approaches but is unlikely to exceed them due to shallow
  adapter, multi-task Stage-2 training, and absent domain-aligned
  inductive biases.
  Wild splits invert the comparison: GM-25 gives small latent gain
  (per-player signal not in engine); Low-Time gives no latent gain
  (BT4 carries no clock state---any LLAMIA advantage is LLM-side);
  Elo-Gap gives moderate latent gain.
\end{itemize}

% =============================================================================
\subsection{Move Annotation (Rationale)}
\label{sec:app_landscape_annotation}
% =============================================================================

\paragraph{What the task requires.}
Given a position and the played move, generate a natural-language
explanation.
The headline metric averages the \emph{Planning} and \emph{Comparative}
subcategories---the two that require genuine strategic reasoning---excluding
Description, Quality, and Context which are more surface-level.

\paragraph{What verbalization cannot express.}
The engine's internal state at any position encodes far more than a move
ranking.
Its activations represent learned lookahead---compressed representations
of tactical sequences three to eight moves deep, including threats the
opponent avoids and defensive resources the opponent fails to find.
They encode attack geometry: which diagonals bear pressure on the king,
which squares in the castled structure are weak, how coordinated pieces
create mating threats.
They encode strategic themes: that a pawn recapture on f6 permanently
damages the dark-square complex, that an isolated d-pawn will become an
endgame target, that the bishop pair grows in strength as exchanges open
the position.

Verbalization collapses all of this into three objects: a scalar evaluation
($+0.5$), a ranked move list (Nxd5, Bg5, Re1), and a principal variation
in algebraic notation (Nxd5 exd5 Bxf6 gxf6 Qh5).
The PV shows \emph{what} happens after the move but not \emph{why} it
works---that Nxd5 eliminates the dark-square defender, that gxf6 creates
a permanent structural weakness, that Qh5 exploits the resulting geometry.
The eval says \emph{how much} better the move is but not \emph{what
changed}---whether the advantage stems from king safety, pawn structure,
piece activity, or material.
The same PV (Nxd5 exd5 Bxf6) could arise from exploiting dark-square
weakness, targeting a backward pawn, or simply winning material, and
neither the eval nor the move sequence disambiguates.

This distinction---between what a position \emph{evaluates to} and what
it \emph{means}---is the core of the verbalization debt on annotation.
Both headline subcategories depend on meaning:

\emph{Planning} requires explaining the move's strategic intent.
``This sacrifice opens the g-file and destroys the dark-square shelter''
is a reading of the engine's positional features---king-safety gradients,
pawn-structure assessments, piece-activity patterns.
The verbalized model must \emph{infer} intent from notation; the
internalized model \emph{reads} it from the latent state.

\emph{Comparative} requires explaining what distinguishes the played move
from alternatives.
``After Bg5 instead of Nxd5, the pawn structure remains intact and the
king can castle---the dark-square weakness that makes Nxd5 decisive is
absent.''
This is a reading of the \emph{difference} between two latent states.
Two eval numbers ($+0.5$ vs.\ $+0.2$) confirm one move is better but
carry no information about the positional mechanism behind the gap.

The excluded subcategories---Description (``the bishop moves to g5,
pinning the knight''), Quality (``this move loses a pawn''), Context
(``this is a well-known Sicilian Najdorf line'')---depend on surface
features that verbalization captures adequately.
The verb--latent gap on the headline metric should therefore be
substantially larger than the gap on the full five-subcategory average.

\paragraph{How non-verbalizable information enables collaboration agency.}
The non-verbalizable signal does not merely improve passive perception---it
enables productive collaboration strategies that verbalization structurally
cannot support (\cref{sec:abl_agency_strategies}).

The dominant strategy for Comparative annotation is the \emph{counterfactual
query}: play the alternative move, invoke the agent on the resulting
position, receive fresh latent tokens, and compare the two latent states.
When the model plays Bg5 instead of Nxd5 and compares latent states, the
difference encodes the positional mechanism directly: the dark-square
features are deactivated in the Bg5 line, the king-safety gradient is
flatter, the pawn-structure representation is intact.
This is structurally what expert annotators do: play the alternative on the
board, assess the resulting position, and articulate the contrast.
Under verbalization, the same counterfactual query returns ``eval $+0.2$
instead of $+0.5$''---the model knows the alternative is worse but has no
access to \emph{why}.
The collaboration strategy heatmap (\cref{fig:agency_heatmaps}) confirms
the asymmetry: counterfactual
query is a significant strategy on annotation under internalization, while
engine-follow dominates under verbalization (62--76\%) because querying
alternatives returns the same compressed summary regardless of the question.

For Planning, the relevant strategy is \emph{feature reading}: the model
attends to specific activations in the latent tokens---king-safety neurons,
pawn-structure features, piece-coordination patterns---and grounds its
explanation in what it observes.
Under verbalization, the model has only the eval and PV; it must
\emph{guess} which positional features justify the move, a lossy inference
that is sometimes correct for powerful frontier LLMs but systematically
limited for smaller trained models.

The collaboration agency effect amplifies the perceptual debt: internalization
provides richer information per invocation \emph{and} enables the model to
use that information strategically---invoking the agent on counterfactual
positions, comparing latent states across alternatives, reading different
features for different annotation subcategories.
The verb--latent gap on annotation is therefore large, close to Commentary
rather than to BC\@.
Commentary adds multi-step state tracking across 30+ moves, which pushes
its gap slightly higher, but the per-position depth of annotation---the
need to explain \emph{why} and \emph{what distinguishes}---activates the
same non-verbalizable information and the same collaboration strategies.

\paragraph{Per-system analysis.}

\emph{LLM Only.}
Produces fluent, plausible annotations that are generically true of many
positions but not specific to the given one.
Planning annotations reference common strategic ideas (``develops the bishop
to an active diagonal'') without positional justification; Comparative
annotations mention obvious alternatives (``could have castled'') without
knowing which alternatives exist or their quality.

\emph{Frontier\,+\,Lc0 Verb.}
A powerful frontier LLM can \emph{infer} positional mechanisms from PV
notation: seeing ``Nxd5 exd5 Bxf6 gxf6 Qh5'' it guesses ``dark-square
destruction.''
This inference is often correct, making Frontier competitive---but it
remains inference, not observation, and fails on positions where the
mechanism is non-obvious from the move sequence.
The frontier model's strength is its reasoning capacity; its limitation is
that reasoning over lossy input cannot recover what verbalization discarded.

\emph{Verbalized baselines (SFT and DAPO).}
Training learns efficient extraction of verbal features for annotation.
DAPO teaches the model to make additional tool calls to explore
alternatives, but the verbal returns from each exploration carry limited
information---``eval went from $+0.3$ to $+0.25$ on the alternative line''
does not reveal what changed positionally.
The model becomes skilled at narrating eval numbers without access to the
mechanisms behind them.

\emph{LLAMIA-SFT (latent).}
The latent tokens encode the positional features that justify the move:
the adapter learns to read king-safety activations, pawn-structure
features, and piece-coordination patterns and ground its annotations in
them.
For Comparative, the full policy distribution identifies which alternatives
are genuinely competitive.
However, single-step training does not teach counterfactual exploration---the
model reads the current position's features but does not compare against
alternative positions.

\emph{LLAMIA (latent\,+\,DAPO).}
RL teaches the counterfactual query strategy and sharpens feature reading
for Planning.
The combination of perceptual access (what the latent state contains) and
collaboration agency (how the model uses it) yields a large verb--latent
gap on both headline subcategories---not just Planning but also
Comparative, where counterfactual discrimination is the primary mechanism.

\emph{SCC (dedicated expert).}
Small supervised model trained on 90K games~\cite{zang2019automated}.
Competitive on Quality and Description but limited on Planning and
Comparative due to no lookahead and no access to engine features.

\paragraph{$2{\times}2$ factorial prediction.}
The interaction effect should be strongly positive---close to Commentary.
RL gain under latent substantially exceeds RL gain under verbalization
because the counterfactual query strategy is discovered by RL and
productive only under internalization.
Both headline subcategories contribute: Planning benefits from RL-sharpened
feature reading, Comparative benefits from RL-discovered counterfactual
comparison.
The verb--latent gap is large on both subcategories, not concentrated in
Planning alone.

\paragraph{Expected qualitative signatures.}
Verbalized systems annotate from eval numbers: ``this move improves White's
position from $+0.3$ to $+0.5$, while the alternative only reaches $+0.3$.''
Accurate but shallow---explains the magnitude but not the mechanism.
LLAMIA annotates from positional features: ``Nxd5 eliminates the key
defender of the dark squares; the gxf6 recapture creates permanent
structural damage---compare with Bg5 where Black's pawn structure remains
intact and the king can castle to safety.''
The Planning subcategory is the diagnostic for perceptual access (does the
annotation explain \emph{what the move is trying to achieve}?); the
Comparative subcategory is the diagnostic for counterfactual discrimination
(does the annotation explain \emph{what distinguishes the played move from
its alternatives}?).

\paragraph{G-eval judge bias.}
The G-eval judge (GPT-4o) introduces a systematic confound:
LLM-as-judge evaluators exhibit family preference, scoring outputs from
GPT-family models 5--15~pp higher than equally good non-GPT
outputs~\cite{XXX_panickssery2024}.
GPT-5's outputs are stylistically closer to GPT-4o's prior than
Qwen3-based LLAMIA outputs are; the measured Frontier\,+\,Lc0 G-eval is
therefore likely \emph{inflated} relative to LLAMIA by this bias.
If LLAMIA-14B matches or exceeds Frontier\,+\,Lc0 in raw G-eval despite
this headwind, the underlying quality gap is larger than the numbers
suggest.
Conversely, a per-subcategory decomposition is essential: if Planning and
Comparative favor LLAMIA while Description and Quality favor Frontier, the
aggregate may mask the integration-interface effect.

\paragraph{Predicted behavior summary.}
\begin{itemize}[leftmargin=*, itemsep=2pt, topsep=2pt]
\item \textbf{Scale (4B${\to}$14B):} Large on Planning and Comparative.
  Both require reading positional features from latent tokens and
  grounding them in language; the 14B's reasoning and grounding
  capacity meaningfully exceeds the 4B's.
\item \textbf{SFT vs.\ DAPO:} Substantial under latent, small under
  verbal.
  RL discovers the counterfactual query strategy---play the alternative,
  invoke the agent, compare latent states---productive only when the
  latent tokens encode discriminative positional features.
\item \textbf{Verbalization debt:} Large, close to Commentary on the
  headline subcategories.
  Planning and Comparative depend on \emph{why} and
  \emph{what distinguishes}, which verbalization destroys; the gap on
  the headline metric is substantially larger than on the full
  five-subcategory average.
\item \textbf{vs.\ Frontier\,+\,Lc0:} Competitive on Planning when the
  mechanism is inferable from PV notation (e.g., recognizing
  ``dark-square destruction'' from ``Nxd5\,exd5\,Bxf6\,gxf6''); fails
  when the mechanism is non-obvious from the move sequence.
  G-eval inflated 5--15~pp by GPT-family judge preference.
\item \textbf{vs.\ task experts:} SCC~\cite{zang2019automated}, trained
  on 90K games, is competitive on Quality and Description but limited
  on Planning and Comparative due to no lookahead and no access to
  engine features.
\end{itemize}

% =============================================================================
\subsection{Game Commentary}
\label{sec:app_landscape_commentary}
% =============================================================================

\paragraph{What the task requires.}
The model produces a natural-language commentary across an entire game,
evaluated by G-eval on relevance, completeness, clarity, and fluency.
Commentary is the \emph{maximally multi-step task} in LLAMIA-Bench: the
unit of evaluation is the game, not the position, and narrative coherence
across 30+ moves is required.

\paragraph{Three compounding factors.}
First, \emph{error compounding}: misidentifying the strategic theme at
move~10 (e.g., describing a kingside attack when the real story is a
queenside pawn majority) corrupts all subsequent commentary.
Under verbalization, each tool-call response is an isolated snapshot; the
model must reconstruct the narrative from accumulated text fragments, and
reconstruction errors accumulate.
Under latent internalization, the re-encoded tokens after each state
transition carry the evolving positional character directly.

Second, \emph{state tracking}: the model must track how themes evolve---who
holds the initiative, how the pawn structure transforms, whether a
kingside attack is building or has been defused.
The latent tokens encode the current positional character at each move;
the verbalized eval numbers do not distinguish between a $+0.5$ arising
from structural advantage, tactical initiative, or material compensation.

Third, \emph{counterfactual reasoning}: analysis of LLAMIA's collaboration
strategies shows that counterfactual queries account for the largest share
of collaboration episodes on this task.
The model plays hypothetical continuations, invokes the agent on the
resulting position, and compares latent states to produce explanations
like ``had Black recaptured \ldots{}Nxd5 instead of \ldots{}exd5, the position
remains balanced; after \ldots{}exd5 the structural damage is permanent.''
This pattern is unproductive under verbalization because comparing
``eval $+0.30$'' to ``eval $+0.25$'' carries negligible information about
what positionally changed.

\paragraph{Importance calibration.}
Expert commentary modulates depth by positional significance: critical
moments receive paragraphs of analysis while routine moves get a single
clause.
Three failure modes prevent non-latent systems from reproducing this
distribution.
LLM-Only and Frontier\,+\,Lc0 Verb produce near-uniform explanation
length across positions because neither has a reliable signal for which
moves deserve elaboration; LLAMIA-Verb-DAPO partially corrects this by
flagging large eval swings, but substitutes eval-magnitude pacing for
the substantive significance pacing that human annotators use.
Verbal systems also over-pivot on small eval drifts (0.1--0.3) that
human commentators ignore as ``human moves'' which do not change the
strategic narrative, while missing the distributional cues---sacrifices,
quiet winners, counter-intuitive retreats---that signal an interesting
moment.
The latter is the same non-verbalizable signal that defines the
Interest task; Commentary inherits that debt at every move and
compounds it across the game.

\paragraph{Per-system analysis.}

\emph{LLM Only.}
``Tour-guide'' commentary: describes what happens move by move without
understanding strategic dynamics.
Can identify opening names from pretraining (reliable for mainline
openings), then degenerates into generic filler: ``an interesting
continuation that creates tension.''
Completely misses critical moments because it cannot assess evaluation
shifts.

\emph{Frontier\,+\,Lc0 Verb.}
The frontier model compensates for the verbal limitation through aggressive
tool use (10+ calls per move) and strong long-context reasoning.
It detects critical moments via large eval swings and generates coherent
narratives from eval trajectories: ``after 23\ldots{}Bd6?, the position
shifts from roughly equal to clearly favoring White.''
GPT-5's chess pretraining knowledge allows it to \emph{infer} positional
mechanisms from eval changes: when it sees a 1.0-eval drop after
\ldots{}gxf6, it can guess ``probably weakened the dark squares.''
This inference is often correct, making the frontier baseline
\emph{surprisingly strong}---potentially surpassing smaller trained verb
models through sheer reasoning capacity.

\emph{LLAMIA-Verb-SFT.}
Trained on single-step (current verbal snapshot, prior commentary context)
$\to$ next segment.
Learns the mapping from verbal features to commentary language but not
narrative planning.
Produces locally coherent segments that may not connect into a global arc.
Commentary is evaluation-driven: tracks numbers rather than positional themes.

\emph{LLAMIA-Verb (DAPO).}
RL over full-game rollouts teaches the model to maintain narrative coherence
and to identify when to elaborate (critical moments) vs.\ summarize
(routine play).
The DAPO reward (G-eval on full-game commentary) directly optimizes for
narrative quality.
This yields the \emph{largest SFT--DAPO gap under verbalization} of any
task, because multi-step narrative planning is exactly what RL teaches.
But the information bottleneck persists: the model becomes a skilled narrator
of evaluation curves without access to the positional mechanisms behind
them.

\emph{LLAMIA-SFT (latent).}
Even under single-step SFT, the latent tokens carry substantially more
information per move than verbalized text.
The model can produce positionally grounded commentary: ``the bishop pair
gains power as the position opens.''
The fresh tokens at each state transition give the model a running
``reading'' of the position's character that verbal systems lack entirely.
We predict that \emph{latent SFT surpasses verbal DAPO on
commentary}: if confirmed, this would demonstrate that on multi-step tasks
the information channel matters more than the optimization procedure.

\emph{LLAMIA (latent\,+\,DAPO).}
The full pipeline: RL teaches counterfactual exploration, narrative pacing,
and critical-moment detection from changes in latent-token patterns (not
just eval swings).
The counterfactual strategy produces commentary that connects causes to
consequences across multiple moves: ``the dark-square weakness from
12\ldots{}gxf6 becomes the defining feature---White's Nf5 exploits this
directly, and the alternative \ldots{}gxf5 creates a passed pawn that
becomes devastating with the bishop pair.''
If RL discovers the counterfactual strategy during training, this
combination of latent access and multi-step optimization should yield the
highest G-eval of any system.

\paragraph{$2{\times}2$ factorial prediction.}
Commentary should exhibit the strongest super-additive interaction of any
task.
RL gain under latent should exceed RL gain under verbal because RL discovers
the counterfactual query strategy, which is productive only when latent
tokens provide discriminative signal between alternative positions.
Latent gain under DAPO should exceed latent gain under SFT because RL
teaches the model to \emph{exploit} the richer signal through multi-step
narrative planning.
Concretely:

\begin{center}
\small
\begin{tabular}{lcc}
\toprule
& \textbf{SFT} & \textbf{DAPO} \\
\midrule
Verbalized & low & moderate \\
Latent     & moderate--high & \textbf{highest} \\
\bottomrule
\end{tabular}
\end{center}

\noindent The interaction term---the extent to which the joint
latent\,+\,DAPO gain exceeds the sum of individual gains---should be
largest on Commentary among all tasks.

\paragraph{Expected qualitative signatures.}
\begin{itemize}[leftmargin=*, itemsep=1pt]
\item \emph{LLM Only}: ``Kasparov plays Nc3, developing another piece.''
  Board description without strategic dynamics.
\item \emph{Verb (SFT)}: ``The eval is $+0.3$ here.
  After Nf5 it becomes $+0.8$.''
  Segment-to-segment incoherent, numbers without mechanism.
\item \emph{Verb (DAPO)}: ``The turning point comes at move~23 where the
  eval swings by 1.2.''
  Correctly identifies turning points but explains them as eval events, not
  positional changes.
\item \emph{Frontier}: ``After gxf6, White seems to have a slight edge,
  probably due to the weakened dark squares.''
  Correct inference from eval + domain knowledge, but hedged because the
  model infers rather than observes.
\item \emph{LLAMIA (SFT)}: ``After the exchange on d5, the pawn structure
  favors White---the isolated d5 pawn is a long-term target.''
  Positionally grounded, reads features from latent tokens, but does not
  explore alternatives.
\item \emph{LLAMIA (DAPO)}: ``Had Black recaptured \ldots{}Nxd5 instead of
  \ldots{}exd5, the position remains balanced---the knight on d5 is
  powerful but White's bishop pair compensates.
  After \ldots{}exd5, the structural damage is permanent.''
  Counterfactual comparison, causal reasoning, narrative connecting past
  decisions to future consequences.
\end{itemize}

\paragraph{Predicted behavior summary.}
\begin{itemize}[leftmargin=*, itemsep=2pt, topsep=2pt]
\item \textbf{Scale (4B${\to}$14B):} Large ($2.7{\times}$ in
  \cref{fig:scaling_bars}); multi-step state tracking and narrative
  pacing across 30+ moves benefit substantially from larger LLMs.
\item \textbf{SFT vs.\ DAPO:} \emph{Largest gap in the benchmark}, under
  both interfaces but especially under latent.
  RL teaches narrative pacing, critical-moment detection from latent
  changes (not just eval swings), and counterfactual exploration.
\item \textbf{Verbalization debt:} Large, slightly above Annotation.
  Per-move debt (positional mechanisms, attack geometry) compounds
  with multi-step state-tracking error; the $2{\times}2$ interaction
  is the \emph{strongest super-additive effect in the benchmark}.
\item \textbf{vs.\ Frontier\,+\,Lc0:} Surprisingly strong via aggressive
  tool use (10+ calls per move) and long-context reasoning over eval
  trajectories.
  Inflated by GPT-family G-eval bias; matching or exceeding Frontier
  in raw G-eval indicates a larger underlying quality gap.
\item \textbf{vs.\ task experts:} No prior model covers full-game
  commentary; SCC and other single-move annotation models address
  isolated positions only.
  The competitive surface is set by Frontier\,+\,Lc0 alone.
\end{itemize}

% =============================================================================
\subsection{Cross-Task Patterns}
\label{sec:app_landscape_cross}
% =============================================================================

Three structural patterns emerge across tasks:

\paragraph{The verb--latent gap scales with non-verbalizable information
and collaboration agency.}
Interest shows the largest relative gap: the engine's distributional
features (policy shape, value gradients) that define interest have no
verbal equivalent, and no collaboration strategy under verbalization can
recover them.
Commentary and Annotation show large gaps driven by both sources of debt.
The engine's lookahead, attack geometry, and strategic themes cannot be
verbalized, and the counterfactual query strategy that exploits these
features is productive only under internalization.
Commentary additionally accumulates debt through multi-step state tracking
across 30+ moves, placing it slightly above Annotation.
BC shows a moderate gap: top-$k$ covers ${\sim}$60--70\% of human moves,
leaving the residual distributional structure---which moves are
``conceptually nearby but tactically inferior'' in the engine's
representation space---accessible only to the latent channel and
load-bearing for skill-conditioned prediction.
Difficulty shows a slightly smaller gap than BC: solution depth is a
strong verbal proxy that absorbs most of the variance, and the latent
residual captures computational-character features (forced-move
sequences, back-rank motifs, distractor sharpness) orthogonal to length.
Both are limited by single-step prediction, which caps the
collaboration-agency contribution.

\paragraph{The SFT--DAPO gap scales with task horizon and collaboration
agency.}
Interest (single-step perceptual judgment): minimal RL benefit.
Difficulty (single-step, benefits from simulated solving): small benefit.
BC (single-step move prediction): small benefit.
Annotation (single-step per position, but counterfactual exploration and
feature reading are load-bearing): substantial benefit---RL discovers the
counterfactual query strategy that is productive only under internalization.
Commentary (full-game narrative, 30+ step horizon): largest benefit---RL
teaches narrative pacing, counterfactual exploration, and critical-moment
detection from latent-token changes.
The interaction between interface and optimization is super-additive
on tasks where RL discovers collaboration strategies (counterfactual
queries, feature reading) that are productive under latent but not under
verbalization---primarily Annotation and Commentary.

\paragraph{Integration interface dominates backbone scale.}
Across the BC literature, domain-aligned representation and conditioning
yield larger gains than parameter scaling: Chessformer-79M surpasses
GPT-3.5-175B by 3.4~pp on move-matching despite being $2000{\times}$
smaller~\cite{XXX_chessformer,zhang2024humanalignedchessbitsearch}.
ALLIE's ablations confirm the model is data-constrained, not
parameter-constrained.
In our setting, the same Qwen3 backbone sits inside three tiers (no
subagent, verbalized, latent); only the Lc0 signal varies.
The predicted ordering---Verb ${<}$ Latent at fixed backbone by a larger
margin than 4B ${<}$ 14B at fixed interface---would extend the literature's
finding to the case where the ``architecture'' is held constant and only the
information channel changes.
For G-eval--based tasks (Annotation, Commentary), the GPT-4o judge's
documented family preference~\cite{XXX_panickssery2024} inflates
Frontier\,+\,Lc0 scores by an estimated 5--15~pp relative to Qwen3-based
LLAMIA outputs, so raw G-eval comparisons between Frontier and LLAMIA
should be interpreted with this confound in mind.
Per-subcategory decomposition and secondary metrics (BLEU-2, perplexity)
provide partial mitigation.
